% GENERATED FILE. Edit canonical lessons, then run npm run generate:books.
% canonical-source-hash: fnv1a64:864d6da8dc281764
% canonical-lessons: PT-C473-gentileza, PT-C473-delicado, PT-C473-grosseiro, PT-C473-antipatico, PT-C473-ofender

\chapter{Gentis e grosseiros}
\label{ch:pt-gentis-e-grosseiros}
\hlchaptermodality{473}

\begin{chapteropening}
\textbf{By the end of this chapter:} \emph{I can describe polite and rude behaviour.}

Complete the last lesson of chapter 473: I can describe polite and rude behaviour.
\end{chapteropening}

\section[a gentileza]{a gentileza — kindness, courtesy}

\label{lesson:PT-C473-gentileza}



\begin{quote}
Before the new one: say the Portuguese for to suspect, then the Portuguese for suspicious.
\end{quote}

\subsection*{You'll want to know: a gentileza}
\textbf{a gentileza} — \textquotedblleft{}kindness, courtesy\textquotedblright{}.

\begin{quote}
Foi uma gentileza do vizinho ajudar-nos com as malas.
\end{quote}

\begin{quote}
\emph{It was kind of the neighbour to help us with the cases.}
\end{quote}

\emph{Por gentileza} — please, kindly (a little formal).

\begin{cousinweb}
\emph{gentileza} is built on \emph{gentil} \textquotedblleft{}kind, courteous\textquotedblright{}, from Latin \emph{gentīlis} \textquotedblleft{}of the same clan\textquotedblright{} (Lewis \& Short), built on \emph{gēns} \textquotedblleft{}a clan, a family\textquotedblright{}.
\end{cousinweb}

\begin{grammarlens}[title={What you've built}]
One more word for telling, explaining and describing at B1.
\end{grammarlens}

\subsection*{Your turn}
\begin{itemize}
\raggedright
  \item \emph{Say it:} \emph{a gentileza}
  \item \emph{Say it:} \emph{a gentileza}, once more
  \item \emph{Say it:} a sentence of your own with \emph{a gentileza}
  \item \emph{Recall:} say the Portuguese for to suspect, then the Portuguese for suspicious, then say \emph{a gentileza} again
\end{itemize}

\begin{tcolorbox}[breakable,colback=teal!4,colframe=teal!35!black,title={Before you move on}]
What does a gentileza mean? (\textquotedblleft{}Kindness, courtesy\textquotedblright{}.) Say it once more.
\end{tcolorbox}
\section[delicado]{delicado — polite, considerate; delicate}

\label{lesson:PT-C473-delicado}



\begin{quote}
Before the new one: say the Portuguese for suspicious, then the Portuguese for kindness.
\end{quote}

\subsection*{You'll want to know: delicado}
\textbf{delicado} — \textquotedblleft{}polite, considerate; delicate\textquotedblright{}.

\begin{quote}
Foi muito delicado da tua parte telefonar.
\end{quote}

\begin{quote}
\emph{It was very thoughtful of you to phone.}
\end{quote}

Of things it means delicate or fragile: \emph{uma situação delicada}.

\begin{cousinweb}
\emph{delicado} is a learned borrowing from Latin \emph{dēlicātus} \textquotedblleft{}charming, dainty, delicate\textquotedblright{} (EtymDB).
\end{cousinweb}

\begin{grammarlens}[title={What you've built}]
One more word for telling, explaining and describing at B1.
\end{grammarlens}

\subsection*{Your turn}
\begin{itemize}
\raggedright
  \item \emph{Say it:} \emph{delicado}
  \item \emph{Say it:} \emph{delicado}, once more
  \item \emph{Say it:} a sentence of your own with \emph{delicado}
  \item \emph{Recall:} say the Portuguese for suspicious, then the Portuguese for kindness, then say \emph{delicado} again
\end{itemize}

\begin{tcolorbox}[breakable,colback=teal!4,colframe=teal!35!black,title={Before you move on}]
What does delicado mean? (\textquotedblleft{}Polite, considerate; delicate\textquotedblright{}.) Say it once more.
\end{tcolorbox}
\section[grosseiro]{grosseiro — rude, coarse}

\label{lesson:PT-C473-grosseiro}



\begin{quote}
Before the new one: say the Portuguese for kindness, then the Portuguese for polite.
\end{quote}

\subsection*{You'll want to know: grosseiro}
\textbf{grosseiro} — \textquotedblleft{}rude, coarse\textquotedblright{}.

\begin{quote}
O empregado foi grosseiro com os clientes.
\end{quote}

\begin{quote}
\emph{The waiter was rude to the customers.}
\end{quote}

Of cloth or work: rough, badly made.

\begin{cousinweb}
\emph{grosseiro} is built on \emph{grosso} \textquotedblleft{}thick, coarse\textquotedblright{}, which is inherited from Late Latin \emph{grossus} \textquotedblleft{}thick\textquotedblright{} (EtymDB).
\end{cousinweb}

\begin{grammarlens}[title={What you've built}]
One more word for telling, explaining and describing at B1.
\end{grammarlens}

\subsection*{Your turn}
\begin{itemize}
\raggedright
  \item \emph{Say it:} \emph{grosseiro}
  \item \emph{Say it:} \emph{grosseiro}, once more
  \item \emph{Say it:} a sentence of your own with \emph{grosseiro}
  \item \emph{Recall:} say the Portuguese for kindness, then the Portuguese for polite, then say \emph{grosseiro} again
\end{itemize}

\begin{tcolorbox}[breakable,colback=teal!4,colframe=teal!35!black,title={Before you move on}]
What does grosseiro mean? (\textquotedblleft{}Rude, coarse\textquotedblright{}.) Say it once more.
\end{tcolorbox}
\section[antipático]{antipático — unfriendly, unpleasant}

\label{lesson:PT-C473-antipatico}



\begin{quote}
Before the new one: say the Portuguese for polite, then the Portuguese for rude.
\end{quote}

\subsection*{You'll want to know: antipático}
\textbf{antipático} — \textquotedblleft{}unfriendly, unpleasant\textquotedblright{}.

\begin{quote}
O novo vizinho é um pouco antipático.
\end{quote}

\begin{quote}
\emph{The new neighbour is a bit unfriendly.}
\end{quote}

The opposite of \emph{simpático}.

\begin{cousinweb}
\emph{antipático} is built on \emph{antipatia} \textquotedblleft{}dislike\textquotedblright{}, a learned borrowing from Greek \emph{antipátheia} \textquotedblleft{}antipathy\textquotedblright{} (EtymDB), built on \emph{anti} \textquotedblleft{}against\textquotedblright{} and \emph{páthos} \textquotedblleft{}feeling\textquotedblright{}.
\end{cousinweb}

\begin{grammarlens}[title={What you've built}]
One more word for telling, explaining and describing at B1.
\end{grammarlens}

\subsection*{Your turn}
\begin{itemize}
\raggedright
  \item \emph{Say it:} \emph{antipático}
  \item \emph{Say it:} \emph{antipático}, once more
  \item \emph{Say it:} a sentence of your own with \emph{antipático}
  \item \emph{Recall:} say the Portuguese for polite, then the Portuguese for rude, then say \emph{antipático} again
\end{itemize}

\begin{tcolorbox}[breakable,colback=teal!4,colframe=teal!35!black,title={Before you move on}]
What does antipático mean? (\textquotedblleft{}Unfriendly, unpleasant\textquotedblright{}.) Say it once more.
\end{tcolorbox}
\section[ofender]{ofender — to offend}

\label{lesson:PT-C473-ofender}



\begin{quote}
Before the new one: say the Portuguese for rude, then the Portuguese for unfriendly.
\end{quote}

\subsection*{You'll want to know: ofender}
\textbf{ofender} — \textquotedblleft{}to offend\textquotedblright{}.

\begin{quote}
As palavras dele ofenderam a minha mãe.
\end{quote}

\begin{quote}
\emph{His words offended my mother.}
\end{quote}

\emph{Ofender-se} — to take offence.

\begin{cousinweb}
\emph{ofender} comes from Latin \emph{offendere} \textquotedblleft{}to strike against, to offend\textquotedblright{} (EtymDB; Lewis \& Short).
\end{cousinweb}

\begin{grammarlens}[title={What you've built}]
One more word for telling, explaining and describing at B1.
\end{grammarlens}

\subsection*{Your turn}
\begin{itemize}
\raggedright
  \item \emph{Say it:} \emph{ofender}
  \item \emph{Say it:} \emph{ofender}, once more
  \item \emph{Say it:} a sentence of your own with \emph{ofender}
  \item \emph{Recall:} say the Portuguese for rude, then the Portuguese for unfriendly, then say \emph{ofender} again
\end{itemize}

\begin{tcolorbox}[breakable,colback=teal!4,colframe=teal!35!black,title={Before you move on}]
What does ofender mean? (\textquotedblleft{}To offend\textquotedblright{}.) Say it once more.
\end{tcolorbox}
